\documentclass[11pt,a4paper]{article}

\usepackage[utf8]{inputenc}
\usepackage[T1]{fontenc}
\usepackage{amsmath,amssymb,amsthm,mathtools}
\usepackage{algorithm}
\usepackage{algpseudocode}
\usepackage{booktabs}
\usepackage[numbers,sort&compress]{natbib}
\usepackage[hidelinks]{hyperref}
\usepackage{cleveref}
\usepackage{microtype}
\setlength{\emergencystretch}{3em}

\theoremstyle{plain}
\newtheorem{theorem}{Theorem}[section]
\newtheorem{proposition}[theorem]{Proposition}
\newtheorem{corollary}[theorem]{Corollary}
\theoremstyle{definition}
\newtheorem{definition}[theorem]{Definition}
\theoremstyle{remark}
\newtheorem{remark}[theorem]{Remark}

\newcommand{\R}{\mathbb{R}}
\newcommand{\N}{\mathbb{N}}
\newcommand{\norm}[1]{\left\lVert #1 \right\rVert}
\newcommand{\abs}[1]{\left\lvert #1 \right\rvert}
\newcommand{\ip}[2]{\left\langle #1, #2 \right\rangle}
\newcommand{\lean}[1]{\path{#1}}
\DeclareMathOperator{\divw}{div}
\DeclareMathOperator{\rank}{rank}
\DeclareMathOperator{\range}{range}
\DeclareMathOperator{\TV}{TV}

\title{The Dynamics of an Evolving Knowledge Graph:\\
Velocity, Flows, Ripples, and the Macro Regime,\\
Proved Before Code}
\author{mind-palace contributors}
\date{November 7, 2026}

\begin{document}
\maketitle

\begin{abstract}
A knowledge graph that is kept over time is a sequence of graphs: notes are written and
rewritten, citations are added and severed, weights drift. We ask what the state, direction and
magnitude of such a graph mean, what a sum over all walks costs, how a flow on the edges is to be
read, and whether a macro regime can be built from micro fluctuations, and we answer each question
with theorems machine-checked in Lean~4 over Mathlib before any code relies on them. The sum over
all walks of a damped relation is a resolvent, one sparse solve, with an entrywise tail bound of
$s^n/(1-s)$ after $n$ terms and a refuted control at spectral radius one. The least-squares
potential of a weighted edge flow is $L_w^{+}\divw_w Y$, its residual divergence-free and its
uniqueness the kernel of the weighted Laplacian (HodgeRank). A function of a symmetric operator
acts on each eigenvector by the function of its eigenvalue, so the heat kernel commutes with the
projection onto any set of modes and the dropped part decays at the rate of the gap below it. The
Laplacian of the quotient by a partition is $P^{\top} L P$, coarse-graining never lowers the
spectral gap, and a block read on its own under-estimates nearness. A spectral gap stated as a
Poincar\'e inequality is exactly a bound under the nonzero eigenvalues, total variation is bounded
by $\sqrt{N}/2$ times the $L^2$ distance, and the heat kernel and the lazy walk contract every
mean-zero field geometrically at the gap's rate. Each theorem is attacked on an exact rational
shadow with a negative control before it is proved. We state plainly which parts of the
statistical-physics analogy are theorems on a finite graph and which are hypotheses that need
data the graph does not yet have.
\end{abstract}

\section{Introduction}

A personal knowledge graph is a small, exact, content-addressed record: a few dozen notes, their
citations, a similarity backbone, and a version chain per note. Its retrieval algebra is already
proved: the walk sum of a killed random walk is well posed and its free energy is a quasi-metric;
the Hodge 1-Laplacian of the flag complex splits every edge flow into a gradient, a curl and a
harmonic part; the heat kernel is a semigroup whose diffusion distance is a metric; transports
along supersessions have known norms. What was missing is time. When a value changes behind a
fixed identity, when weights drift, when a shock is propagated along a relation, what are the
objects, and what can be known?

This article answers six questions in order, each from a machine-checked theorem.

\begin{enumerate}
\item \emph{The objects.} An identity is a box; the value it holds is re-recorded at each cut;
  its velocity over a window of cuts is a $0$-cochain per pair of cuts. A field among the edges
  is either a flow (a $1$-cochain, antisymmetric) or a change of weights (symmetric), and the two
  have different readings.
\item \emph{Path sums.} The sum over all walks of a damped relation is the resolvent
  $(1 - zA)^{-1}$: one sparse solve, never an enumeration. We prove the convergence condition in
  its operator-norm form, an entrywise tail bound, and the failure at spectral radius one.
\item \emph{Flows.} A weighted edge flow has a least-squares potential $L_w^{+}\divw_w Y$, whose
  residual is divergence-free and which is unique up to the kernel of $L_w$; the potential ranks
  the nodes and the residual's energy measures inconsistency.
\item \emph{Macro from micro.} The low modes of the Laplacian are the macro variables: the heat
  kernel evolves them exactly and the dropped modes decay at the rate of the gap below them.
  Coarse-graining by a partition is exact on block-constant fields, never lowers the gap, and a
  block read on its own under-estimates nearness.
\item \emph{Mixing.} A spectral gap as a Poincar\'e inequality is exactly a bound under the
  nonzero eigenvalues; total variation is bounded by $\sqrt{N}/2$ times the $L^2$ distance; the
  heat kernel and the lazy walk contract mean-zero fields geometrically.
\item \emph{What can be known.} The propagation of a known shock is exact; the shocks are not
  predicted by anything here; the regime of the graph's own evolution is a hypothesis.
\end{enumerate}

\paragraph{Contributions.}
(i) A model of a value's versions, edge flows, the Neumann partial sums, the ripple of a shock,
the modes of a symmetric operator and the quotient of a graph by a partition, in Lean~4 over
Mathlib (\lean{Palace/Dynamics/Defs.lean}). (ii) Six families of theorems, DYN-1 to DYN-6
(\lean{Palace/Dynamics.lean}), each stated below with a proof sketch naming its declaration, and
replayed by Lean's kernel and five independent checkers. (iii) Property attacks on exact rational
shadows with negative controls (\lean{PalaceTest/Dynamics.lean}, seeds 330 to 344). (iv) A plain
account of which parts of the fluid and statistical-physics analogy are theorems on a finite graph
and which are hypotheses, and of the sorts and verbs a query language needs to express the
dynamics, with the order to build them.

\section{Background and related work}

\paragraph{Randomized shortest paths and the walk sum.}
The retrieval kernel of the graph is the randomized-shortest-path (RSP) family of
\citet{saerens2009rsp} and \citet{kivimaki2014rsp}: a random walk killed at rate $1-z$ per step,
biased by a Boltzmann factor at inverse temperature $\beta$, whose fundamental matrix
$Z = (I - W_\beta)^{-1}$ sums over all walks and whose free energy
$\varphi_{ij} = -\beta^{-1}\log(z_{ij}/z_{jj})$ interpolates between the commute-time distance
and the shortest-path distance. The fundamental matrix of an absorbing chain is classical
\citep{kemeny1960finite}. The convergence of the Neumann series $\sum_k M^k$ when a norm of $M$ is
below one, and the sharpness of the spectral-radius condition through Gelfand's formula, are in
\citet[Chapter~5]{horn2013matrix}.

\paragraph{Hodge theory on graphs and ranking.}
The Hodge decomposition of edge flows on a simplicial complex, $C^1 = \range\partial_1^{\top}
\oplus \range\partial_2 \oplus \ker L_1$, and its reading of a pairwise comparison as a global
ranking plus local and global inconsistencies, is HodgeRank \citep{jiang2011hodgerank}; the
operator-theoretic account is \citet{lim2020hodge}. The same decomposition reads a game as a
potential game plus a harmonic game \citep{candogan2011flows}.

\paragraph{Spectral graph theory, diffusion and mixing.}
The Laplacian's spectrum, the heat kernel $e^{-tL}$ and the Rayleigh quotient are the subject of
\citet{chung1997spectral}; diffusion maps read the low modes of a diffusion as coordinates
\citep{coifman2006diffusion}. The spectral gap as a Poincar\'e inequality between a Dirichlet
form and a variance, and the geometric bounds it gives on mixing, are in
\citet{diaconis1991geometric} and \citet{levin2017mixing}. The three mixing statements of
Section~\ref{sec:mixing} follow the shapes of the switch-chain formalization in
\citet{openai2026switch}: a gap as a Poincar\'e inequality, total variation against $L^2$ by
$\sqrt{N}/2$, and geometric $L^2$ decay per step; they are restated here over the modes of any
weighted graph's Laplacian, and nothing is copied.

\paragraph{Formalization.}
The proofs are in Lean~4 \citep{demoura2021lean4} over Mathlib \citep{mathlib2020}, at the pin
v4.36.0-rc2. Mathlib supplies the geometric series in a complete normed ring, the spectral theorem
for Hermitian matrices with its orthonormal eigenvector basis and continuous functional calculus,
and the row-sum operator norm on matrices. It has, at this pin, no spectral gap, no Cheeger
inequality, no heat kernel and no Markov mixing, so Section~\ref{sec:mixing} is ours.

\section{Theory}
\label{sec:theory}

Throughout, $V$ is a finite set of nodes, $E$ a finite set of oriented edges, matrices are real,
and $x \cdot y$ is the dot product. Every theorem below is proved in \lean{Palace/Dynamics.lean};
the declaration is named after each statement.

\subsection{The objects}

\begin{definition}[velocity]
Let $x : \N \to V \to \mathcal{E}$ assign to each cut $c$ and identity $v$ the value $x_c(v)$ alive
at $c$, in an additive group $\mathcal{E}$ (the reals for a number, $\R^d$ for an embedding). The
velocity over the window $c_0 \le c_1$ is the $0$-cochain
$v_{c_0 c_1}(i) = x_{c_1}(i) - x_{c_0}(i)$ (\lean{velocity}).
\end{definition}

\begin{theorem}[DYN-1]
\label{thm:velocity}
Velocity adds along the chain of cuts, $v_{ab} + v_{bc} = v_{ac}$ (\lean{velocity_add}), and
the displacement over a window is the sum of its steps,
$v_{a,a+n} = \sum_{k<n} v_{a+k,a+k+1}$ (\lean{velocity_eq_sum}). For real-valued $x$, the
coboundary of the velocity is the change of the coboundary,
$(dv_{ab})(u,w) = (dx_b)(u,w) - (dx_a)(u,w)$ (\lean{coboundary_velocity}), and on any complex
with $\partial_1\partial_2 = 0$ the edge field $\partial_1^{\top} v_{ab}$ is orthogonal to every
curl $\partial_2 y$ (\lean{grad_velocity_dot_curl}).
\end{theorem}

\begin{proof}[Proof sketch]
The first two are the group identities, the second through Mathlib's telescoping sum
\lean{Finset.sum_range_sub}. The third is pointwise. The fourth is the orthogonality of gradients
and curls in the Hodge family (\lean{Hodge.grad_dot_curl}), applied to $v_{ab}$.
\end{proof}

A flow is a $1$-cochain $Y : E \to \R$, antisymmetric under reversal; a change of weights is a
function on unordered edges, symmetric. \Cref{thm:velocity} says the edge field induced by a
changing potential is a gradient, so circulation, when there is any, does not come from a changing
potential.

\subsection{The Neumann resolvent}

\begin{definition}
For $M \in \R^{V\times V}$ the Neumann partial sum is $S_n(M) = \sum_{k<n} M^k$ (\lean{neumann});
for a relation $A$, a damping $z$ and a shock $\delta : V \to \R$, the ripple is
$(1 - zA)^{-1}\delta$ (\lean{ripple}). $M$ is row-bounded by $s$ when $M \ge 0$ entrywise and
every row sums to at most $s$.
\end{definition}

\begin{theorem}[DYN-2, the series]
\label{thm:neumann}
In a complete normed ring, when $\norm{x} < 1$ the powers of $x$ are summable and
$(1-x)\sum_k x^k = (\sum_k x^k)(1-x) = 1$ (\lean{neumann_tsum}). For a real matrix under the
row-sum operator norm, when $\norm{M} < 1$ the powers are summable entrywise and
$(1-M)^{-1} = \sum_k M^k$ (\lean{inv_one_sub_eq_tsum}). A matrix row-bounded by $s \ge 0$ has
$\norm{M} \le s$ (\lean{linfty_opNorm_le_of_rows}), so when $s<1$ its walk sum is its inverse
(\lean{inv_one_sub_eq_tsum_of_rows}); in particular the fundamental matrix of the killed
walk is the sum over all killed walks (\lean{fundamental_eq_tsum}) and the ripple of a shock is
$\sum_k z^k A^k \delta$ (\lean{ripple_eq_tsum}).
\end{theorem}

\begin{proof}[Proof sketch]
The ring statement is Mathlib's \lean{geom_series_mul_neg} and \lean{mul_neg_geom_series}.
For matrices, the row-sum norm \lean{Matrix.linftyOpNormedRing} is a local instance; the
matrices are complete under it because its uniform structure is the product's, which Lean checks
by unfolding (\lean{inferInstanceAs}); then \lean{Matrix.nonsing_inv_eq_ringInverse} and
Mathlib's \lean{geom_series_eq_inverse} give the inverse. The norm bound is
\lean{Matrix.linfty_opNorm_def}, the supremum of the row sums of absolute values, with the
entries nonnegative.
\end{proof}

\begin{theorem}[DYN-2, the tail]
\label{thm:tail}
Let $M$ be row-bounded by $s<1$. Then the rows of $M^n$ sum to at most $s^n$
(\lean{sum_pow_le}), the rows of $(1-M)^{-1}$ sum to at most $1/(1-s)$ (\lean{sum_inv_le}),
$(1-M)^{-1} - S_n(M) = (1-M)^{-1} M^n$ (\lean{inv_sub_neumann}), and entrywise
\[
0 \;\le\; (1-M)^{-1}_{ij} - S_n(M)_{ij} \;\le\; \frac{s^n}{1-s}
\]
(\lean{neumann_tail}).
\end{theorem}

\begin{proof}[Proof sketch]
The row sums of powers by induction on $n$ with $M^{n+1} = M^n M$. The row sums of the inverse by
the extremal principle of the retrieval family: $Z = 1 + MZ$ entrywise
(\lean{Retrieval.inv_apply_eq}), so at the row $m$ where the row sum $r$ is largest,
$r_m \le 1 + s\,r_m$. The identity is $(1-M)S_n = 1 - M^n$ (Mathlib's \lean{mul_neg_geom_sum})
multiplied by $(1-M)^{-1}$, which exists by RSP-1's \lean{det_one_sub_ne_zero}. The bound
sums $Z_{ik}(M^n)_{kj} \le Z_{ik} s^n$ over $k$, using that an entry of a nonnegative matrix is
at most its row sum.
\end{proof}

\begin{theorem}[DYN-2, the control]
\label{thm:control}
If every row of $M$ sums to exactly $1$ then $\det(1-M) = 0$
(\lean{det_one_sub_eq_zero_of_rows_eq_one}) and the partial sums $S_n(M)$ converge to no
matrix (\lean{neumann_not_tendsto}).
\end{theorem}

\begin{proof}[Proof sketch]
$(1-M)\mathbf{1} = 0$ with $\mathbf{1} \ne 0$ (\lean{Matrix.exists_mulVec_eq_zero_iff}). Each row
of $M^k$ sums to $1$, so each row of $S_n(M)$ sums to $n$; a limit would make the sequence $n$
converge in $\R$, against \lean{tendsto_natCast_atTop_atTop}.
\end{proof}

\begin{theorem}[DYN-5, path independence across a transport]
\label{thm:intertwine}
If $PM = M'P$ and both $1-M$ and $1-M'$ are invertible then
$P(1-M)^{-1} = (1-M')^{-1}P$ (\lean{resolvent_intertwine}). For an injective $\sigma : A \to B$
with indicator $P_\sigma$, if $M'(\sigma a, \sigma a') = M(a,a')$ for all $a, a'$ and
$M'(b, \sigma a) = 0$ for every $b$ outside the range of $\sigma$, then
$P_\sigma^{\top} M = M' P_\sigma^{\top}$ (\lean{transpose_indicator_mul}), and the ripple of a
shock pushed forward along $\sigma$ is the ripple of the pushed shock
(\lean{push_ripple}).
\end{theorem}

\begin{proof}[Proof sketch]
From $PM = M'P$, $(1-M')P = P(1-M)$; multiply on the left by $(1-M')^{-1}$ and on the right by
$(1-M)^{-1}$. The entry $(b,a)$ of $P_\sigma^{\top}M$ is $M(a_0,a)$ when $b = \sigma a_0$ and $0$
otherwise, by injectivity; that of $M'P_\sigma^{\top}$ is $M'(b,\sigma a)$.
\end{proof}

\subsection{HodgeRank: the least-squares potential}

\begin{definition}
For an incidence matrix $B \in \R^{V\times E}$ and weights $w : E \to \R_{\ge 0}$, let
$\ip{a}{b}_w = \sum_e w_e a_e b_e$ (\lean{wdotE}), $L_w = B\,\mathrm{diag}(w)\,B^{\top}$
(\lean{weightedLap}), $\divw_w Y = B(w \circ Y)$ (\lean{divergence}), the potential
$s = L_w^{+}\divw_w Y$ (\lean{potential}) with $L_w^{+}$ the Moore--Penrose inverse of the
Spectral family, and the residual $r = Y - B^{\top}s$ (\lean{residual}).
\end{definition}

\begin{theorem}[DYN-3]
\label{thm:hodgerank}
For $w \ge 0$: $x\cdot L_w x = \ip{B^{\top}x}{B^{\top}x}_w$ (\lean{dotProduct_weightedLap});
$L_w k = 0$ iff $w \circ B^{\top}k = 0$ (\lean{weightedLap_mulVec_eq_zero_iff}); the residual
is divergence-free, $\divw_w r = 0$ (\lean{divergence_residual}); for every $x$,
$\ip{r}{r}_w \le \ip{Y - B^{\top}x}{Y - B^{\top}x}_w$ (\lean{wdotE_residual_le}), with equality
iff $L_w(s - x) = 0$ (\lean{wdotE_residual_eq_iff}). On the flag complex of a connected graph
with every edge weight positive, $L_w k = 0$ implies $k$ constant
(\lean{eq_of_weightedLap_mulVec_eq_zero}).
\end{theorem}

\begin{proof}[Proof sketch]
The energy identity moves $B$ across the dot product. The kernel: a zero energy forces each
$w_e (B^{\top}k)_e^2 = 0$. The divergence is orthogonal to the kernel of $L_w$, since
$\divw_w Y \cdot k = Y \cdot (w \circ B^{\top}k)$, so $L_w L_w^{+}$ fixes it by the Penrose
equations (\lean{mulVec_pinv_mulVec_of_orthogonal}, Rayleigh family), which is the normal
equation $L_w s = \divw_w Y$. Pythagoras: $Y - B^{\top}x = r + B^{\top}(s-x)$ and $r$ is
$w$-orthogonal to every gradient (\lean{wdotE_residual_grad}), so the energies add
(\lean{wdotE_sub_grad}). On the flag complex, $(d_1^{\top}k)_e = k(e_2) - k(e_1)$ and positive
weights force equality across every edge; a walk carries it across the component.
\end{proof}

\subsection{The modes of a symmetric operator}

Let $M$ be real symmetric with Mathlib's eigenvalues $\lambda_j$ and orthonormal eigenvectors
$\varphi_j$ (\lean{mode}); let $P_S f = \sum_{j\in S}\ip{f}{\varphi_j}\varphi_j$ (\lean{modeProj})
and $H_t = e^{-tM}$ the heat kernel, defined through the continuous functional calculus
(\lean{heat}, Heat family).

\begin{theorem}[DYN-4]
\label{thm:modes}
For any $f : \R \to \R$, $f(M)\varphi_j = f(\lambda_j)\varphi_j$ (\lean{cfc_mulVec_mode}) and
$x\cdot f(M)x = \sum_j f(\lambda_j)\ip{x}{\varphi_j}^2$ (\lean{dotProduct_cfc_mulVec}).
The heat kernel commutes with the projection onto any set of modes and acts on them by their
eigenvalues,
$H_t P_S f = \sum_{j\in S} e^{-t\lambda_j}\ip{f}{\varphi_j}\varphi_j = P_S H_t f$
(\lean{heat_modeProj}). For $t \ge 0$ and $\mu \le \lambda_j$ for all $j \notin S$,
\[
\norm{H_t(f - P_S f)}^2 \le e^{-2\mu t}\,\norm{f - P_S f}^2
\]
(\lean{heat_sub_modeProj_le}).
When $\mu \le \lambda_j$ for every $j$ with $\ip{x}{\varphi_j} \ne 0$,
$\mu\norm{x}^2 \le x\cdot Mx$ (\lean{mul_dotProduct_self_le_of_modes}).
\end{theorem}

\begin{proof}[Proof sketch]
The spectral theorem writes $f(M) = U\,\mathrm{diag}(f\circ\lambda)\,U^{*}$ with $U$ the
eigenvector unitary; $U^{*}\varphi_j = e_j$ and $Ue_j = \varphi_j$
(\lean{star_eigenvectorUnitary_mulVec}, \lean{eigenvectorUnitary_mulVec}). Parseval,
$x\cdot y = \sum_j\ip{x}{\varphi_j}\ip{y}{\varphi_j}$, is the orthonormal basis's
\lean{sum_inner_mul_inner}; with $y = f(M)x$ and the symmetry of $f(M)$ the quadratic form
follows. The commutation expands $P_S f$ and uses linearity; the decay writes
$\norm{H_t g}^2 = g\cdot H_{2t}\,g$ by the semigroup law (\lean{heat_add}) and bounds each
$e^{-2t\lambda_j}$ by $e^{-2t\mu}$ on the modes $g$ lives on, which for $g = f - P_S f$ are the
dropped ones (\lean{sub_modeProj_dotProduct_mode}). The last statement bounds each term of the
energy sum.
\end{proof}

\subsection{Coarse-graining by a partition}

\begin{definition}
For a partition $\pi : V \to B$ let $P_{vb} = [\pi(v) = b]$ (\lean{indicator}) and
$A_\pi = P^{\top}AP$ (\lean{quot}). The pull-back $\pi^{*}h = h\circ\pi$ and the push-forward
$\pi_{*}g$, summing over each block, are the Transport family's.
\end{definition}

\begin{theorem}[DYN-5]
\label{thm:coarse}
Over any commutative ring, $P^{\top}\,L(A)\,P = L(A_\pi)$ (\lean{lap_quot}), so
$\pi^{*}h\cdot L(A)\,\pi^{*}h = h\cdot L(A_\pi)\,h$ (\lean{dotProduct_lap_pull}). For symmetric
$A \ge 0$ with connected support and $\mu$ at most every nonzero eigenvalue of $L(A)$: a field $x$
with $x\cdot\mathbf{1} = 0$ has $\mu\norm{x}^2 \le x\cdot L(A)x$ (\lean{gap_le_dotProduct_lap}),
and a block field $h$ with $\pi^{*}h\cdot\mathbf{1} = 0$ has
$\mu\sum_b \abs{b}\,h_b^2 \le h\cdot L(A_\pi)h$ (\lean{gap_le_dotProduct_lap_quot}). A block of
a partition, read with the graph's degrees, has nearness at most the graph's
(\lean{nearnessOn_block_le}, RSP-5 on the block).
\end{theorem}

\begin{proof}[Proof sketch]
$P^{\top}\mathrm{diag}(d)P = \mathrm{diag}(P^{\top}d)$ (\lean{transpose_indicator_mul_diagonal})
and the degree of $A_\pi$ is $P^{\top}(A(P\mathbf{1})) = P^{\top}d$ since $P\mathbf{1} = \mathbf{1}$
(\lean{indicator_mulVec_one}). A mean-zero field has no weight on a zero mode of a connected
support, because such a mode is a kernel vector, hence constant
(\lean{eq_of_lap_mulVec_eq_zero}, from Rayleigh's \lean{orthogonal_ker_lap}), so
\Cref{thm:modes}'s last statement applies (\lean{dotProduct_mode_eq_zero}); the block form is
the energy identity with $\norm{\pi^{*}h}^2 = \sum_b\abs{b}h_b^2$ (TRN-2,
\lean{pull_dotProduct_self}).
\end{proof}

\subsection{Mixing}
\label{sec:mixing}

\begin{definition}
$A$ has spectral gap at least $\gamma$ when $\gamma\norm{f}^2 \le f\cdot L(A)f$ for every $f$ with
$f\cdot\mathbf{1} = 0$ (\lean{HasSpectralGapAtLeast}).
\end{definition}

\begin{theorem}[DYN-6]
\label{thm:mixing}
Let $A$ be symmetric with nonnegative entries.
\begin{enumerate}
\item On a connected support, $\gamma$ at most every nonzero eigenvalue is a gap
  (\lean{hasSpectralGapAtLeast_of_eigenvalues}); conversely a gap $\gamma$ is at most every
  nonzero eigenvalue (\lean{le_eigenvalue_of_hasSpectralGapAtLeast}); and a gap bounds the
  Dirichlet energy of any $f$ below by $\gamma$ times its variance about its mean
  (\lean{HasSpectralGapAtLeast.variance}).
\item For every $x : V \to \R$, $\tfrac12\sum_i\abs{x_i} \le \tfrac{\sqrt{N}}{2}\norm{x}$
  with $N = \abs{V}$ (\lean{half_sum_abs_le}).
\item On a connected support with $\gamma$ at most every nonzero eigenvalue, for
  $x\cdot\mathbf{1} = 0$ and $t\ge 0$: $\norm{H_t x}^2 \le e^{-2\gamma t}\norm{x}^2$
  (\lean{heat_energy_le_of_gap}) and
  $\tfrac12\sum_i\abs{(H_t x)_i} \le \tfrac{\sqrt N}{2}\,e^{-\gamma t}\norm{x}$
  (\lean{half_sum_abs_heat_le}); and for $\varepsilon \ge 0$ with
  $\varepsilon\lambda_j \le 1$ on every eigenvalue,
  $\norm{(1-\varepsilon L)^n x}^2 \le (1-\varepsilon\gamma)^{2n}\norm{x}^2$
  (\lean{lazy_pow_energy_le}).
\end{enumerate}
\end{theorem}

\begin{proof}[Proof sketch]
(1) The forward direction is \Cref{thm:coarse}. For the converse, a mode with a nonzero eigenvalue
is orthogonal to $\mathbf{1}$, since $\lambda_j\ip{\varphi_j}{\mathbf{1}} = (L\varphi_j)\cdot\mathbf{1}
= \varphi_j\cdot L\mathbf{1} = 0$ (\lean{mode_dotProduct_one}), and the Poincar\'e inequality on
$\varphi_j$ reads $\gamma \le \lambda_j$. The variance form subtracts the mean, which $L$ kills.
(2) Cauchy--Schwarz, Mathlib's \lean{sq_sum_le_card_mul_sum_sq}. (3) \Cref{thm:modes}'s
decay with the mean-zero field's modes; the total-variation form composes with (2); the lazy walk
acts on $\varphi_j$ by $1 - \varepsilon\lambda_j$ (\lean{one_sub_smul_mulVec_mode}), and a
symmetric operator acting on each mode by $g_j$ has
$\norm{N^n x}^2 = \sum_j g_j^{2n}\ip{x}{\varphi_j}^2$ (\lean{dotProduct_self_mulVec_of_modes},
\lean{pow_mulVec_mode}), bounded on the modes $x$ lives on by $\rho^{2n}$ when
$\abs{g_j} \le \rho$ there (\lean{pow_energy_le}).
\end{proof}

\section{Concepts and intuition}

\subsection{The resolvent is the path sum}

A path sum asks for every walk from $i$ to $j$, each weighted by the product of its steps and
damped by $z$ per step. \Cref{thm:neumann} says that on a finite graph this is a geometric
series of matrices and that its sum is an inverse: the whole sum is one sparse linear solve
$(1-zA)Z = 1$, polynomial in the number of nodes, never an enumeration of walks, whose number is
exponential in their length. When an iterative answer is wanted, \Cref{thm:tail} gives the
accuracy of $n$ terms and says the partial sums never overshoot. \Cref{alg:ripple} is the
computation, with the tail bound as its accuracy line. The killed walk of the retrieval kernel
is the case the graph already computes, and the free energy $-\beta^{-1}\log(z_{ij}/z_{jj})$ is a
free energy in the statistical-physics sense: $z_{ij}$ is a partition function over walks at
inverse temperature $\beta$, positive, and subadditive along paths. The control
(\Cref{thm:control}) is the unkilled walk: at spectral radius one there is no resolvent and no sum.

\begin{algorithm}[t]
\caption{Ripple of a shock $\delta$ along a relation $A$ with damping $z$.}
\label{alg:ripple}
\begin{algorithmic}[1]
\Require $A \ge 0$ with rows summing to at most $1$; $0 \le z < 1$; a shock $\delta$; a tolerance
  $\epsilon$ or \textsc{exact}
\State $M \gets zA$; \quad $s \gets \max_i \sum_j M_{ij}$ \Comment{$s \le z < 1$}
\If{\textsc{exact}}
  \State \Return the solution $r$ of $(1 - M)\,r = \delta$ \Comment{one sparse solve}
\Else
  \State $n \gets \lceil \log(\epsilon(1-s)) / \log s \rceil$ \Comment{\Cref{thm:tail}}
  \State $r \gets \delta$; \quad $u \gets \delta$
  \For{$k = 1, \dots, n-1$}
    \State $u \gets M u$; \quad $r \gets r + u$ \Comment{$r = S_k(M)\delta$, below the sum}
  \EndFor
  \State \Return $r$ \Comment{entrywise within $\epsilon\norm{\delta}_\infty$ of the sum}
\EndIf
\end{algorithmic}
\end{algorithm}

Across time, two orders are possible: propagate the shock on the graph at one cut and carry the
result forward along the supersession, or carry the shock forward and propagate it at the next
cut. \Cref{thm:intertwine} says they agree exactly when the transport intertwines the relations,
which holds for a merge-free transport that carries every weight forward with no new node feeding
the carried ones. The chain-map condition of the superconnection (every citation carries forward
or collapses) is necessary and not sufficient: a carried citation with a changed weight breaks the
intertwining where the curvature is zero. Across a merge the orders differ (the control at seed
340).

\subsection{The potential is the ranking}

A flow on the edges, a lead-lag estimate say, is read by its Hodge decomposition: a gradient
$\partial_1^{\top}s$, a curl, and a harmonic part. \Cref{thm:hodgerank} says the potential
$s = L_w^{+}\divw_w Y$ is the least-squares ranking of the nodes: $s(j) - s(i)$ is the part of
the flow across $i \to j$ that a global order explains, and the residual's weighted energy
$\ip{r}{r}_w$ is the inconsistency, zero exactly when the flow is a gradient. \Cref{alg:hodgerank}
is the computation. The curl is disagreement that closes within a filled triangle; the harmonic
part is disagreement that closes only around a hole, of dimension the first Betti number of the
complex. A whole that is more than its parts, in this reading, is the harmonic part: circulation
that no node potential and no local triangle accounts for.

\begin{algorithm}[t]
\caption{HodgeRank: the potential and the inconsistency of a weighted flow.}
\label{alg:hodgerank}
\begin{algorithmic}[1]
\Require incidence $B$ ($+1$ at an edge's head, $-1$ at its tail), weights $w \ge 0$, a flow $Y$
\State $L_w \gets B\,\mathrm{diag}(w)\,B^{\top}$; \quad $d \gets B(w\circ Y)$
\State $s \gets L_w^{+} d$ \Comment{\Cref{thm:hodgerank}: the normal equations $L_w s = d$}
\State $r \gets Y - B^{\top}s$ \Comment{divergence-free: $B(w\circ r) = 0$}
\State \Return $s$ up to a constant per component, and $\ip{r}{r}_w$
\end{algorithmic}
\end{algorithm}

The projections do not commute with a mask on the nodes: a filled triangle with the unit
circulation is pure curl, but restricted to two of its vertices it is one edge with flow one, a
pure gradient. A mask changes the complex, and the decomposition belongs to the complex. This is
derived with the witness and not yet machine-checked.

\subsection{Modes are the macro variables}

A field on the nodes is a micro state. Its component along the constants is the macro state, the
mean, conserved by the diffusion since $L\mathbf{1} = 0$; its components along the other modes
are the fluctuations. \Cref{thm:modes} says the heat equation evolves each mode by its own rate
with no coupling, so the coefficients on any set of modes are exact macro variables, and that
what is dropped decays at least as fast as $e^{-\mu t}$ for $\mu$ the gap below the kept modes.
\Cref{thm:coarse} says a coarse-graining by blocks is exact on block-constant fields and can lose
slow modes but never create one: the quotient graph's Fiedler quotient is at least the fine gap.
\Cref{thm:mixing} says the fluctuations average out at a rate set by one number of the graph, the
gap, in $L^2$ and in total variation, in continuous time and for the lazy walk. The analogy with
a diffusing fluid is exact at every $N$ because the heat equation on a graph is a diffusion; what
a small $N$ loses is not the truth of the theorems but the smoothness of a continuum: a graph of a
few dozen nodes has a handful of modes, each legible, none of them a density.

\section{Evaluation}

\paragraph{What is machine-checked.}
The six families are proved in Lean~4 at Mathlib pin v4.36.0-rc2 and pass the repository's gate
(\lean{formal/check.sh}): the library builds with warnings fatal; the attacks run; Batteries'
environment linters pass; the sorry count is held at zero; no axiom beyond \lean{propext},
\lean{Classical.choice} and \lean{Quot.sound} is used (3853 declarations audited on the run that
included this family); the export of every declaration and everything it uses, 45768
declarations from \lean{Init} up, is replayed through Lean's kernel and five independent checkers
(leanchecker-paranoid, lean4lean, nanoda, con-leche, con-ron), all of which accept. The 248
ledger statements judged by the comparator are unchanged by this family, which adds design
theorems only.

\paragraph{What is attacked.}
Each attackable theorem is run against 500 random instances of an exact rational shadow before
it is proved, with a negative control, a nearby false statement that must be refuted
(\Cref{tab:attacks}). The heat kernel and the modes are exponentials and eigenvectors, which do
not run, so DYN-4 and DYN-6 are proved and not attacked.

\begin{table}[t]
\centering
\caption{Attacks (seed, 500 instances each, none discarded) and their controls.}
\label{tab:attacks}
\small
\begin{tabular}{@{}p{3.4cm}p{3.6cm}p{5.0cm}@{}}
\toprule
attack & shadow & control (seed, refuted within) \\
\midrule
tail in $[0, s^n/(1-s)]$, $n \le 3$ (330) & $5\times5$ matrices with rows summing under $1$,
  Gauss--Jordan & without the factor $1/(1-s)$ (331, 14); signed entries overshoot (332, 6); rows
  summing to $1$ make $1-M$ singular (333, first instance) \\
HodgeRank on $K_5$ (334) & weights in $\{0,\tfrac14,\tfrac12,\tfrac34,1\}$, $L_w^{+}$ by the
  Spectral shadow & the potential of the unweighted divergence is not divergence-free (335, 14) \\
$P^{\top}LP = L_\pi$ (336) and its energy form (338) & $5\times5$ integer matrices, three blocks &
  lifted back to the fine graph it is not a Laplacian (337, 10) \\
ripple across a permutation (339) & products of transpositions & across a merge the ripple is not
  path-independent (340, 7) \\
velocity telescoping (341) & values from lists & the sum of the steps' sizes is not the
  displacement's (342, 26) \\
RSP-5 on a block (343) & the killed walk at $\beta=1$, $z=\tfrac12$ & the block's own degrees
  invent nearness (344, 31) \\
\bottomrule
\end{tabular}
\end{table}

\paragraph{What is measured.}
Read-only on the live instance (2026-11-07): the versions store holds 29 identities and 29
versions, every version at sequence number one, all recorded within one second of the last rebuild
(2026-10-28). There is one generation, so there is no displacement to measure and no series: the
velocity of \Cref{thm:velocity} is a definition with its laws, waiting on the first supersession.
The similarity backbone measured on 2026-11-04 has 38 notes and 38 edges, and the relation graph
20 edges on 12 notes, so a spectrum here is a short list of eigenvalues, each legible.

\section{Limitations}

The sharp convergence condition of the walk sum, a spectral radius below one, is proved in its
operator-norm form and cited in its spectral form: Mathlib at the pin proves Gelfand's formula
over $\mathbb{C}$ and the real-matrix statement is not machine-checked. The non-commutation of
the Hodge projections with a mask is derived with a witness and not yet machine-checked. The
regime of the graph's own evolution, whether the sequence of snapshots has a stationary law, a
temperature, a fluctuation-dissipation balance, is a hypothesis: every theorem here concerns
diffusion and propagation on a graph at a cut, and nothing here gives the edits that produce the
next cut a law. The shocks are acts of authorship outside the graph; the dynamics propagate them
exactly and do not forecast them. A thermodynamic limit does not exist for a record that is
finite and exact by design.

\section{Future work}

The query language of the graph has sorts for fields, kernels, metrics, relations, chains and
transports, and lacks the objects of this article. The order to build, each unit transcribing
its theorems and property-testing the transcription against the Lean statements:
(1) a Flow sort (a $1$-cochain with the edge weights as its inner product, a vector space and not
a cone), with the verbs \texttt{gradient}, \texttt{curl}, \texttt{harmonic},
\texttt{potential} and \texttt{inconsistency}, licensed by the Hodge decomposition and
\Cref{thm:hodgerank}, with the mask witness as a refused rewrite;
(2) \texttt{ripple $\langle$rel$\rangle$ from $\langle$anchor$\rangle$ [damping $z$]}, the
resolvent of a directed relation at a shock, licensed by \Cref{thm:neumann,thm:tail}, with the
tail bound as the trace's accuracy line and the symmetric \texttt{diffuse} as its special case;
(3) the adjectives \texttt{at scale $t$} (the heat kernel's column at the anchor, licensed by the
semigroup law and \Cref{thm:modes,thm:mixing}) and \texttt{coarse by $\langle$partition$\rangle$}
(evaluation on the quotient with the fine degrees, licensed by \Cref{thm:coarse} and its
own-degrees control);
(4) \texttt{velocity of $\langle$identity$\rangle$} over a window of cuts, licensed by
\Cref{thm:velocity}, typed as a rate with its clock declared, waiting on the version moves at the
node grain, the snapshot at a past cut, and a second generation in the store. Machine-checking the
mask witness and the spectral form of the convergence condition are the next theorems.

\section{Conclusion}

An evolving knowledge graph has exact objects: a value's velocity over a window of cuts, flows on
its edges with a least-squares ranking and an inconsistency, the ripple of a shock as a resolvent
computed by one solve, and the low modes of its Laplacian as macro variables that evolve exactly
and whose complement averages out at the rate of the spectral gap. Each is a theorem checked by
six kernels before any code relies on it, and each theorem names its hypotheses. What remains is
not a theorem: the shocks themselves, and the regime of the graph's own evolution, which the
graph will be able to test once it has history.

\bibliographystyle{plainnat}
\bibliography{refs}

\end{document}
